% Las siete particiones de 5; etiquetas separadas de los diagramas.
\begin{tikzpicture}[x=.29cm,y=.29cm]
\foreach \cols/\pos/\lab/\paridad in {
{1,1,1,1,1}/0/{[1,1,1,1,1]}/{+1},
{2,1,1,1}/7/{[2,1,1,1]}/{-1},
{2,2,1}/14/{[2,2,1]}/{+1},
{3,1,1}/21/{[3,1,1]}/{+1},
{3,2}/28/{[3,2]}/{-1},
{4,1}/35/{[4,1]}/{-1},
{5}/42/{[5]}/{+1}}
{
  \pgfmathsetmacro{\anchodiagrama}{max(\cols)}
  \node[anchor=south,font=\scriptsize] at (\pos+2,1) {$\lab$};
  \foreach \longitud [count=\fila from 0] in \cols {
    \foreach \col in {1,...,\longitud} {
      \draw[black] (\pos+2-\anchodiagrama/2+\col-1,-\fila) rectangle ++(1,-1);
    }
  }
  \node[anchor=north,font=\small] at (\pos+2,-6) {$\paridad$};
}
\end{tikzpicture}
